\chapter{Dinamica delle Rotazioni}

\section{Introduzione}

In questo capitolo studiamo il comportamento dei sistemi in rotazione, introducendo le tre grandezze vettoriali che ne governano la dinamica: il \textbf{momento di un vettore}, il \textbf{momento di una forza} e il \textbf{momento angolare}.

Tutte e tre hanno una caratteristica in comune che conviene mettere subito in chiaro: non sono proprietà intrinseche del corpo, ma dipendono dal \emph{punto} rispetto al quale le si calcola. Prima ancora di definirle occorre dunque fissare tale punto di riferimento.

\begin{definizione}{Polo della Rotazione}
Si chiama \textbf{polo della rotazione} il punto $O$ rispetto al quale vengono calcolati tutti i momenti. Per comodità lo si sceglie in genere coincidente con l'origine del sistema di riferimento cartesiano adottato.
\end{definizione}

\begin{attenzione}{Ogni Momento Va Riferito a un Polo}
Parlare del ``momento angolare di un corpo'' senza specificare il polo è privo di significato: lo stesso corpo, nello stesso istante, ha momenti angolari diversi rispetto a poli diversi. Per questo la notazione porta sempre il polo a pedice, $\vec{L}_O$ e $\vec{\tau}_O$.
\end{attenzione}

\section{Momento di un Vettore}

Per quantificare l'attitudine di una grandezza vettoriale a produrre una rotazione attorno a un punto si definisce il suo momento. Dato un generico vettore $\vec{B}$ applicato nel punto $A$, il suo momento rispetto al polo $O$ è il prodotto vettoriale fra il raggio vettore $\vec{r} = \vec{OA}$ e il vettore stesso:
\begin{equation}
    \vec{M}_B = \vec{OA} \times \vec{B}
\end{equation}

La retta su cui giace $\vec{B}$ prende il nome di \textbf{retta d'azione}. Il momento gode di una proprietà notevole: dipende soltanto dalla retta d'azione, non dal particolare punto di applicazione scelto lungo di essa. Possiamo cioè far scorrere il vettore lungo la propria retta d'azione senza alterarne l'effetto rotazionale.

\begin{osservazione}{Indipendenza dal Punto di Applicazione}
Scegliamo come punto di applicazione un altro punto $A'$ della stessa retta d'azione e calcoliamo il momento corrispondente:
\begin{equation}
    \vec{M}_B' = \vec{OA}' \times \vec{B}
    = (\vec{OA} + \vec{AA}') \times \vec{B}
    = \vec{OA} \times \vec{B} + \vec{AA}' \times \vec{B}
\end{equation}
Poiché $A$ e $A'$ appartengono entrambi alla retta d'azione di $\vec{B}$, lo spostamento $\vec{AA}'$ è parallelo a $\vec{B}$; il prodotto vettoriale fra vettori paralleli è nullo, perché l'angolo compreso vale zero e $\sin 0 = 0$. Il secondo termine si annulla e resta
\begin{equation}
    \vec{M}_B' = \vec{OA} \times \vec{B} = \vec{M}_B
\end{equation}
\end{osservazione}

La conseguenza pratica è che il modulo del momento si calcola comodamente come
\begin{equation}
    |\vec{M}_B| = |\vec{B}| \, b
\end{equation}
dove $b = |\vec{OH}|$ è il \textbf{braccio}, cioè la distanza fra il polo e la retta d'azione, misurata perpendicolarmente. Tutta l'informazione geometrica rilevante è contenuta in quella sola distanza.

\begin{figure}[ht]
    \centering
    \begin{tikzpicture}[x={(-0.55cm,-0.38cm)}, y={(1cm,0cm)}, z={(0cm,1cm)},
                        >=stealth, scale=1.15, font=\small]
        \coordinate (O) at (0,0,0);

        % Piano xy, per dare profondità
        \fill[mypen, opacity=0.05] (-1,-3,0) -- (3.2,-3,0) -- (3.2,4.3,0) -- (-1,4.3,0) -- cycle;

        % Assi
        \draw[->, mypen!55, thick] (0,0,0) -- (3.2,0,0) node[anchor=north east] {$x$};
        \draw[->, mypen!55, thick] (0,0,0) -- (0,4.3,0) node[anchor=north west] {$y$};
        \draw[->, mypen!55, thick] (0,0,0) -- (0,0,4.3) node[anchor=south] {$z$};
        \node[anchor=north west, mypen, xshift=2pt] at (O) {$O$};

        % Retta d'azione
        \draw[dashed, thick, mypen!70] (2,-3,0) -- (2,4.6,0)
            node[above, anchor=south west, mypen] {retta d'azione};

        % Braccio
        \draw[very thick, boxatt] (O) -- (2,0,0) node[pos=0.55, anchor=south, boxatt, yshift=2pt] {$b$};
        \fill[boxatt] (2,0,0) circle (1.6pt) node[anchor=north, boxatt, yshift=-2pt] {$H$};

        % Raggi vettori
        \coordinate (A)  at (2,-2.2,0);
        \coordinate (Ap) at (2,2.3,0);
        \draw[->, thick, boxese] (O) -- (A)  node[midway, anchor=south east, boxese] {$\vec{r}$};
        \draw[->, thick, boxese!70, dashed] (O) -- (Ap) node[pos=0.72, anchor=south, boxese, yshift=2pt] {$\vec{r}\,'$};

        % Vettori B, applicati in A e in A'
        \fill[mypen] (A) circle (1.6pt) node[anchor=east, mypen, xshift=-2pt] {$A$};
        \fill[mypen] (Ap) circle (1.6pt) node[anchor=north, mypen, yshift=-2pt] {$A'$};
        \draw[->, ultra thick, boxoss] (A) -- (2,-0.85,0) node[midway, anchor=north, boxoss, yshift=-1pt] {$\vec{B}$};
        \draw[->, thick, boxoss!45, dashed] (Ap) -- (2,4.1,0) node[midway, anchor=north west, boxoss] {$\vec{B}$};

        % Momento
        \draw[->, ultra thick, notered] (O) -- (0,0,3.5)
            node[anchor=west, notered] {$\vec{M}_B = \vec{r} \times \vec{B}$};

        % Simboli di ortogonalità
        \draw[thin, mypen] (0,0,0.3) -- (0.35,0,0.3) -- (0.35,0,0);
        \draw[thin, mypen] (1.75,0,0) -- (1.75,0.25,0) -- (2.0,0.25,0);
    \end{tikzpicture}
    \caption{Momento del vettore $\vec{B}$ rispetto al polo $O$. I punti $A$ e $A'$ sono due scelte arbitrarie del punto di applicazione lungo la stessa retta d'azione: il momento $\vec{M}_B$, ortogonale al piano individuato da $\vec{r}$ e $\vec{B}$, è il medesimo in entrambi i casi e il suo modulo vale $|\vec{B}|\,b$, dove $b$ è il braccio.}
    \label{fig:momento_vettore}
\end{figure}

\section{Grandezze della Dinamica Rotazionale}

Applicando la definizione appena data ai due vettori fondamentali della dinamica --- la forza e la quantità di moto --- si ottengono le due grandezze su cui si regge l'intero capitolo.

\begin{definizione}{Momento di una Forza}
Se il vettore considerato è una forza $\vec{F}$, il suo momento rispetto al polo $O$ si chiama \textbf{momento della forza} (o \emph{momento torcente}) e si indica con $\vec{M}_F$ o, più comunemente, con $\vec{\tau}$:
\begin{equation}
    \vec{\tau} = \vec{OA} \times \vec{F} = \vec{r} \times (m\vec{a})
\end{equation}
dove $\vec{r}$ è il raggio vettore del punto di applicazione. Il modulo è massimo quando la forza è perpendicolare al raggio vettore, e in tal caso vale semplicemente
\begin{equation*}
    \tau = m\,r\,a \, \sin\!\left(\frac{\pi}{2}\right) = m\,r\,a
\end{equation*}
\end{definizione}

\begin{definizione}{Momento Angolare}
Se il vettore considerato è la quantità di moto $\vec{p} = m\vec{v}$ di un punto materiale, il suo momento rispetto al polo $O$ si chiama \textbf{momento angolare} e si indica con $\vec{M}_p$ o, più comunemente, con $\vec{L}$:
\begin{equation}
    \vec{L} = \vec{OA} \times \vec{p} = \vec{r} \times (m\vec{v})
\end{equation}
\end{definizione}

Dalla definizione discende immediatamente una proprietà geometrica di grande portata. I vettori $\vec{r}$ e $\vec{v}$ individuano un piano, detto \textbf{piano istantaneo della traiettoria}, e il momento angolare, essendo un prodotto vettoriale, è per costruzione perpendicolare a tale piano.

\begin{osservazione}{Momento Angolare Costante Implica Moto Piano}
Se $\vec{L}$ mantiene direzione e verso costanti, allora anche il piano individuato da $\vec{r}$ e $\vec{v}$ resta fisso nello spazio, e la traiettoria è necessariamente \textbf{piana}. È il motivo per cui i pianeti, soggetti a una forza che conserva il momento angolare, percorrono orbite contenute in un piano anziché curve sghembe nello spazio.
\end{osservazione}

\subsection{Significato Fisico del Momento Angolare}

Per legare il momento angolare allo stato di rotazione conviene esprimere i vettori in coordinate polari, usando il versore radiale $\hat{r}$ e il versore trasverso $\hat{\phi}$:
\begin{itemize}
    \item il raggio vettore ha solo componente radiale, $\vec{r} = r\hat{r}$;
    \item la velocità, essendone la derivata temporale, contiene sia la variazione della distanza sia quella dell'angolo: $\vec{v} = \dot{\vec{r}} = \dot{r}\hat{r} + r\dot{\phi}\hat{\phi}$.
\end{itemize}

Sostituendo nella definizione di momento angolare:
\begin{equation}
    \begin{split}
        \vec{L} &= \vec{r} \times m\vec{v}
        = m \left[ r\hat{r} \times \big(\dot{r}\hat{r} + r\dot{\phi}\hat{\phi}\big) \right] \\
        &= m \left( r\dot{r}\,(\hat{r} \times \hat{r}) + r^2 \dot{\phi}\,(\hat{r} \times \hat{\phi}) \right)
    \end{split}
\end{equation}
Il primo termine si annulla perché $\hat{r} \times \hat{r} = \vec{0}$; il secondo si semplifica ricordando che $\hat{r} \times \hat{\phi} = \hat{z}$, versore uscente dal piano del moto. Indicando con $\omega = \dot{\phi}$ la velocità angolare si ottiene
\begin{equation}
    \vec{L} = m r^2 \dot{\phi} \, \hat{z}
\end{equation}

\begin{osservazione}{Il Momento Angolare Misura la Rotazione, non la Traslazione}
Il risultato è molto istruttivo: $\vec{L}$ si annulla se e solo se $\dot{\phi} = \omega = 0$, cioè quando il moto è puramente radiale. La componente $\dot{r}$ della velocità --- per quanto grande --- non contribuisce affatto. Il momento angolare è dunque una grandezza \textbf{associata esclusivamente alla rotazione} attorno al polo, e il suo valore dipende solo da $r$ e da $\dot{\phi}$.
\end{osservazione}

Il caso notevole è quello in cui si impedisce ogni variazione radiale, $\dot{r} = 0$: il raggio è fissato e il moto è \textbf{circolare}. All'interno di questa famiglia si distinguono due sottocasi:
\begin{itemize}
    \item se l'accelerazione angolare è non nulla ($a_\phi \neq 0$), e dunque $\dot{\phi}$ varia nel tempo, si ha un \textbf{moto circolare vario};
    \item se l'accelerazione angolare è nulla ($a_\phi = 0$), e dunque $\dot{\phi}$ è costante, si ha un \textbf{moto circolare uniforme}.
\end{itemize}

\section{Moto non Planare rispetto a \texorpdfstring{$O$}{O}}

Non sempre il polo $O$ appartiene al piano in cui avviene il moto. Quando il moto circolare si svolge su un piano che non contiene il polo, il momento angolare rispetto a $O$ acquista una struttura più ricca, e conviene calcolarlo scomponendo il raggio vettore.

L'idea è semplice: si proietta il polo $O$ sul piano dell'orbita, ottenendo il centro $O'$ della circonferenza, e si scrive $\vec{r}$ come somma di due contributi ortogonali --- il raggio $\vec{R}$ della circonferenza, che giace nel piano del moto, e la quota $z$ che separa i due poli lungo l'asse di rotazione.

\begin{figure}[ht]
    \centering
    \tdplotsetmaincoords{70}{120}
    \begin{tikzpicture}[tdplot_main_coords, scale=1.45, >=stealth,
                        line cap=round, line join=round, font=\small]
        \def\R{1.5}
        \def\H{2.5}
        \def\angP{45}

        % Assi
        \draw[->, thick, mypen!55] (0,0,0) -- (2.5,0,0) node[anchor=north east] {$x$};
        \draw[->, thick, mypen!55] (0,0,0) -- (0,2.5,0) node[anchor=north west] {$y$};
        \draw[->, thick, mypen!55] (0,0,0) -- (0,0,\H+1.4) node[anchor=south] {$z$};
        \node[anchor=north east, mypen] at (0,0,0) {$O$};

        % Circonferenza e centro proiettato
        \coordinate (Op) at (0,0,\H);
        \node[anchor=north east, mypen] at (Op) {$O'$};
        \fill[mypen] (Op) circle (1.4pt);
        \draw[thick, mypen!75] plot[domain=0:360, samples=60, smooth]
            ({\R*cos(\x)}, {\R*sin(\x)}, \H);

        % Quota z fra i due poli
        \draw[very thick, boxatt] (0,0,0) -- (Op);
        \node[anchor=east, boxatt] at (0,0,{\H/2}) {$z$};

        % Punto P
        \coordinate (P) at ({\R*cos(\angP)}, {\R*sin(\angP)}, \H);
        \fill[notered] (P) circle (1.7pt) node[anchor=south west, mypen] {$P$};

        % Vettori
        \draw[->, thick, boxese] (0,0,0) -- (P)
            node[midway, anchor=north west, boxese] {$\vec{r}$};
        \draw[->, thick, boxoss] (Op) -- (P)
            node[midway, anchor=south, boxoss] {$\vec{R}$};
        \coordinate (V) at ({\R*cos(\angP) - 1.2*sin(\angP)}, {\R*sin(\angP) + 1.2*cos(\angP)}, \H);
        \draw[->, thick, notered] (P) -- (V) node[anchor=west, notered] {$\vec{v}$};

        % Velocità angolare
        \draw[->, thick, mypen] plot[domain=150:-30, samples=30, smooth]
            ({0.42*cos(\x)}, {0.42*sin(\x)}, {\H+0.85});
        \node[mypen] at (0.95, 0.85, {\H+0.95}) {$\vec{\omega}$};

        % Versori locali, sul lato opposto della circonferenza
        \def\angL{170}
        \coordinate (L) at ({\R*cos(\angL)}, {\R*sin(\angL)}, \H);
        \draw[->, very thick, boxese]
            (L) -- ({(\R+0.75)*cos(\angL)}, {(\R+0.75)*sin(\angL)}, \H)
            node[anchor=north east, boxese] {$\hat{\rho}$};
        \draw[->, very thick, boxoss]
            (L) -- ({\R*cos(\angL) - 0.75*sin(\angL)}, {\R*sin(\angL) + 0.75*cos(\angL)}, \H)
            node[anchor=north west, boxoss] {$\hat{\phi}$};
    \end{tikzpicture}
    \caption{Moto circolare su un piano che non contiene il polo $O$. Il raggio vettore si scompone nel raggio $\vec{R} = R\hat{\rho}$ della circonferenza e nella quota $z$ che separa $O$ dal centro proiettato $O'$.}
    \label{fig:moto_non_planare}
\end{figure}

Fissato il piano del moto perpendicolarmente all'asse $\hat{z}$, il centro della circonferenza è $O'$ e la quota che lo separa da $O$ vale $z$. Adottiamo i versori $\hat{\rho}$ (radiale, uscente dal centro $O'$ verso il punto) e $\hat{\phi}$ (trasverso), osservando che il raggio della circonferenza è per costruzione parallelo al primo, $\vec{R} \parallel \hat{\rho}$.

Il raggio vettore rispetto a $O$ si scrive allora come somma dei due contributi ortogonali, $\vec{r} = \vec{R} + \vec{OO}'$, cioè
\begin{equation}
    \vec{r} = R\hat{\rho} + z\hat{z}
\end{equation}
mentre la velocità, essendo il moto circolare, è puramente tangenziale:
\begin{equation}
    \vec{v} = R\omega\hat{\phi}
\end{equation}

Sostituiamo nella definizione di momento angolare e distribuiamo il prodotto vettoriale:
\begin{align}
    \vec{L}_O &= m \, \vec{r} \times \vec{v} = m \, (R \hat{\rho} + z \hat{z}) \times (R \omega \hat{\phi}) \\
              &= m \big[ R \hat{\rho} \times R \omega \hat{\phi} + z \hat{z} \times R \omega \hat{\phi} \big] \\
              &= m \big[ R^2 \omega \, (\hat{\rho} \times \hat{\phi}) + R \omega z \, (\hat{z} \times \hat{\phi}) \big]
\end{align}

Ricordando le regole di moltiplicazione della terna destrorsa, $\hat{\rho} \times \hat{\phi} = \hat{z}$ e $\hat{z} \times \hat{\phi} = -\hat{\rho}$, si ottiene il risultato finale:
\begin{equation}
    \vec{L}_O = m R^2 \omega \, \hat{z} - m R \omega z \, \hat{\rho}
\end{equation}

\begin{osservazione}{Le Due Componenti del Momento Angolare}
Il risultato si legge come somma di due contributi distinti:
\begin{itemize}[leftmargin=1.4em]
    \item $\vec{L}_{O'} = m R^2 \omega \, \hat{z}$, diretto lungo l'asse di rotazione: è esattamente il momento angolare che si avrebbe scegliendo come polo il centro $O'$ della circonferenza, cioè in assenza di sfasamento fra i due poli;
    \item $-\, m R \omega z \, \hat{\rho}$, diretto radialmente verso l'asse: compare unicamente a causa della quota $z$ ed è proporzionale ad essa.
\end{itemize}
Di conseguenza $\vec{L}_O$ \textbf{non è parallelo a} $\vec{\omega}$, e per di più il termine radiale ruota insieme al punto: il momento angolare rispetto a $O$ non è costante in direzione, pur essendolo in modulo. Solo scegliendo il polo sul piano del moto ($z = 0$) le due grandezze tornano parallele.
\end{osservazione}

\section{Forze Centrali}

Si dicono \textbf{centrali} le forze dirette lungo il raggio vettore $\vec{r}$ che congiunge la particella a un punto fisso $O$, detto \textbf{centro di forza}.

\begin{definizione}{Forza Centrale}
Una forza è centrale se la sua retta d'azione passa in ogni istante per il centro di forza $O$, e se il suo modulo dipende soltanto dalla distanza da esso:
\begin{equation}
    \vec{F}_c = \vec{F}(r) = F(r)\,\hat{r}
\end{equation}
dove $F(r)$ è una funzione scalare della sola coordinata radiale $r$. Il segno di $F(r)$ distingue le forze attrattive ($F < 0$) da quelle repulsive ($F > 0$).
\end{definizione}

Le forze fondamentali della natura sono in larga parte di questo tipo: lo sono la forza \textbf{gravitazionale}, la forza \textbf{coulombiana} e la forza \textbf{elastica} isotropa. Lo studio che segue si applica dunque tanto al moto dei pianeti quanto a quello degli elettroni attorno al nucleo.

\subsection{Proprietà Fondamentali}

Le forze centrali possiedono due proprietà generali, dalle quali discende tutto il resto.

\paragraph{1. Sono forze conservative.}
Il lavoro compiuto da una forza centrale dipende unicamente dalle posizioni iniziale e finale. Per dimostrarlo scomponiamo lo spostamento elementare $d\vec{l}$ in una componente radiale $d\vec{r}$ e una trasversale $d\vec{s}$ (lungo il versore $\hat{t}$) e calcoliamo il lavoro elementare:
\begin{equation}
    d\mathcal{L}_{F_c} = \vec{F}_c \cdot d\vec{l}
    = \big(F(r)\hat{r}\big) \cdot \big(dr\,\hat{r} + ds\,\hat{t}\big) = F(r)\,dr
\end{equation}
Il contributo trasversale si annulla perché $\hat{r} \cdot \hat{t} = 0$. Integrando fra due posizioni $\vec{x}_1$ e $\vec{x}_2$:
\begin{equation}
    \mathcal{L}_{F_c} = \int_{\vec{x}_1}^{\vec{x}_2} \vec{F}_c \cdot d\vec{l}
    = \int_{r_1}^{r_2} F(r) \, dr
\end{equation}
L'integrale dipende solo dagli estremi radiali $r_1$ e $r_2$, e non dal cammino percorso per congiungerli: la forza è quindi conservativa e ammette un'energia potenziale.

\begin{figure}[ht]
    \centering
    \begin{tikzpicture}[>=stealth, scale=1.35, font=\small]
        \coordinate (O) at (0,0);
        \coordinate (P) at (2.5,2.0);

        % Traiettoria
        \draw[thick, mypen!45, dashed] (0.8,0.5) to[out=30, in=240] (P) to[out=60, in=190] (4.6,3.5);
        \node[mypen!70, right] at (4.6,3.5) {traiettoria};

        % Direzione radiale
        \draw[dashed, thin, mypen!55] (O) -- (4.1,3.28);

        % Forza centrale, verso O
        \draw[->, ultra thick, boxatt] (P) -- ($(P)!0.55!(O)$)
            node[midway, anchor=south east, boxatt] {$\vec{F}_c$};

        % Punti
        \fill[mypen] (O) circle (1.7pt) node[anchor=north east, mypen] {$O$};
        \fill[notered] (P) circle (1.7pt) node[anchor=north west, mypen] {$P$};

        % Scomposizione dello spostamento
        \coordinate (Pn) at ($(P) + (0.85,0.85)$);
        \coordinate (Pr) at ($(P) + (0.85,0.68)$);
        \draw[->, very thick, notered] (P) -- (Pn)
            node[anchor=south east, notered, yshift=2pt] {$d\vec{l}$};
        \draw[->, very thick, boxese] (P) -- (Pr)
            node[anchor=north west, boxese, xshift=-2pt] {$d\vec{r}$};
        \draw[->, very thick, boxoss] (Pr) -- (Pn)
            node[midway, anchor=west, boxoss] {$d\vec{s}$};
        \draw[dashed, thin, mypen!45] (Pr) -- (Pn);
    \end{tikzpicture}
    \caption{Scomposizione dello spostamento elementare $d\vec{l}$ in componente radiale $d\vec{r}$ e trasversale $d\vec{s}$. La forza centrale $\vec{F}_c$, parallela a $\hat{r}$, compie lavoro soltanto sulla prima.}
    \label{fig:lavoro_forza_centrale}
\end{figure}

\paragraph{2. Il loro momento rispetto a $O$ è nullo.}
Scegliendo come polo il centro di forza stesso, la verifica è immediata:
\begin{equation}
    \vec{\tau}_O = \vec{r} \times \vec{F}_c
    = (r \hat{r}) \times \big(F(r) \hat{r}\big)
    = r F(r) \, (\hat{r} \times \hat{r}) = \vec{0}
\end{equation}
perché $\vec{r}$ e $\vec{F}_c$ sono paralleli per definizione.

\begin{formulario}{Le Due Costanti del Moto in un Campo Centrale}
Da queste due proprietà discendono immediatamente due integrali primi del moto:
\begin{minipage}[t]{0.47\textwidth}
\centering \textbf{Energia totale}
\begin{equation*}
    E = K + U(r) = \text{costante}
\end{equation*}
\centering \footnotesize (la forza è conservativa)
\end{minipage}
\hfill
\begin{minipage}[t]{0.47\textwidth}
\centering \textbf{Momento angolare}
\begin{equation*}
    \frac{d\vec{L}}{dt} = \vec{\tau}_O = \vec{0} \implies \vec{L} = \text{costante}
\end{equation*}
\centering \footnotesize (purché $O \equiv$ centro di forza)
\end{minipage}
\vspace{4pt}
\end{formulario}

Sono queste due costanti a rendere risolubile il problema del moto in un campo centrale, e sarà da esse che ricaveremo, nella sezione conclusiva, la forma generale della traiettoria.

\subsection{Significato Geometrico del Momento Angolare}

La costanza di $\vec{L}$ ammette una lettura puramente geometrica, che storicamente precede di quasi un secolo la formulazione newtoniana. Mentre il punto si muove, il raggio vettore spazza una superficie; calcoliamone il ritmo di crescita.

\begin{figure}[ht]
    \centering
    \begin{tikzpicture}[>=stealth, scale=1.65, font=\small]
        \coordinate (O) at (0,0);
        \coordinate (P1) at (20:3);
        \coordinate (P2) at (50:3.2);

        % Settore spazzato
        \fill[boxatt!15] (O) -- (P1) to[bend right=10] (P2) -- cycle;

        % Traiettoria
        \draw[thick, mypen!45, dashed]
            (5:2.8) to[out=100, in=260] (P1) to[out=80, in=230] (P2) to[out=50, in=200] (66:3.5);
        \node[mypen!70, right] at (66:3.5) {traiettoria};

        % Raggi vettori
        \draw[->, very thick, boxese] (O) -- (P1)
            node[midway, anchor=north west, boxese] {$\vec{r}(t)$};
        \draw[->, very thick, boxese] (O) -- (P2)
            node[midway, anchor=south east, boxese] {$\vec{r}(t+dt)$};

        % Spostamento
        \draw[->, thick, notered] (P1) -- (P2)
            node[midway, anchor=west, notered, xshift=2pt] {$d\vec{r}$};

        % Angolo
        \draw[thick, mypen] (20:0.85) arc (20:50:0.85);
        \node[mypen] at (35:1.1) {$d\varphi$};

        % Punti
        \fill[mypen] (P1) circle (1.6pt) node[anchor=west, mypen, xshift=2pt] {$P(t)$};
        \fill[mypen] (P2) circle (1.6pt) node[anchor=south, mypen, yshift=2pt] {$P(t+dt)$};
        \fill[mypen] (O) circle (1.6pt) node[anchor=north east, mypen] {$O$};

        % Area
        \node[boxatt, font=\bfseries] at (35:1.85) {$dA$};
    \end{tikzpicture}
    \caption{Interpretazione geometrica della costanza del momento angolare: nell'intervallo $dt$ il raggio vettore spazza l'area infinitesima $dA$, assimilabile a un settore circolare di ampiezza $d\varphi$.}
    \label{fig:velocita_areolare}
\end{figure}

Nell'intervallo $dt$ il raggio vettore ruota di $d\varphi$ e spazza un triangolo mistilineo, assimilabile per $d\varphi \to 0$ a un settore circolare di area
\begin{equation}
    dA = \frac{r^2 \, d\varphi}{2}
\end{equation}
Dividendo per $dt$ si ottiene la \textbf{velocità areolare}, e moltiplicando e dividendo per la massa vi si riconosce il momento angolare:
\begin{equation}
    \frac{dA}{dt} = \frac{r^2}{2}\,\frac{d\varphi}{dt}
    = \frac{r^2 \dot{\varphi}}{2} \cdot \frac{m}{m}
    = \frac{|\vec{L}_O|}{2m} = \text{costante}
\end{equation}
avendo usato $|\vec{L}_O| = m r^2 \dot{\varphi}$. Poiché in un campo centrale $|\vec{L}_O|$ è costante, lo è anche la velocità areolare.

\begin{teorema}{Legge delle Aree (Seconda Legge di Keplero)}
In un campo di forze centrali il raggio vettore $\vec{r}$ spazza \textbf{aree uguali in tempi uguali}.
\end{teorema}

\begin{osservazione}{Conseguenza sulla Velocità Tangenziale}
Dalla costanza di $r^2\dot{\varphi}$ segue direttamente
\begin{equation}
    r^2 \dot{\varphi} = (r \dot{\varphi})\, r = v_{\text{tang}} \cdot r = \text{costante}
    \qquad \Longrightarrow \qquad
    v_{\text{tang}} \propto \frac{1}{r}
\end{equation}
La componente tangenziale della velocità decresce dunque in proporzione inversa alla distanza: un pianeta corre più veloce al perielio e più lento all'afelio, esattamente come richiede la legge delle aree.
\end{osservazione}

La costanza della velocità areolare permette infine di legare il \textbf{periodo} orbitale all'area racchiusa dall'orbita:
\begin{equation}
    A_{\text{orb}} = \int_0^T \frac{dA}{dt} \, dt
    = \frac{dA}{dt} \int_0^T dt
    = \frac{|\vec{L}_O|}{2m} \cdot T
    \qquad \Longrightarrow \qquad
    T = \frac{2 \, m \, A_{\text{orb}}}{|\vec{L}_O|}
\end{equation}

\subsection{Energia Potenziale di una Forza \texorpdfstring{$\vec{F}_c \propto r^{-2}$}{Fc proporzionale a r alla meno due}}

Applichiamo quanto visto al caso di gran lunga più importante, quello in cui l'intensità della forza decade come l'inverso del quadrato della distanza --- la legge comune alla gravitazione e all'elettrostatica:
\begin{equation}
    F(r) = \frac{\alpha}{r^2}
\end{equation}
dove la costante $\alpha$ è negativa per una forza attrattiva e positiva per una repulsiva.

Poiché la forza è conservativa, l'energia potenziale si ottiene integrando a partire da una posizione di riferimento $r_0$:
\begin{equation}
    \begin{split}
        U(r) - U(r_0) &= -\,\mathcal{L}_{r_0 \to r} = - \int_{r_0}^{r} \frac{\alpha}{r'^2} \, dr' \\
        &= \int_{r_0}^{r} \left(-\frac{\alpha}{r'^2}\right) dr'
        = \left[ \frac{\alpha}{r'} \right]_{r_0}^{r}
        = \frac{\alpha}{r} - \frac{\alpha}{r_0}
    \end{split}
\end{equation}

Resta da fissare il riferimento. La scelta naturale è porlo all'infinito: a distanza molto grande l'interazione si estingue e sia la forza sia l'energia potenziale tendono a zero. Ponendo dunque $r_0 \to +\infty$ il termine $\alpha/r_0$ si annulla e si ottiene l'espressione definitiva:
\begin{equation}
    U(r) - \cancelto{0}{U(+\infty)} = \frac{\alpha}{r}
    \qquad \Longrightarrow \qquad
    U(r) = \frac{\alpha}{r}
\end{equation}

\begin{attenzione}{La Scelta del Riferimento è Convenzionale}
Il valore assoluto di $U$ non ha significato fisico: solo le \emph{differenze} di energia potenziale sono misurabili, perché è da esse che si ricava il lavoro. Porre $U(+\infty) = 0$ è una convenzione comoda, non un risultato: cambiandola si aggiungerebbe a $U(r)$ una costante additiva senza alterare né le forze né il moto.
\end{attenzione}

\section{Seconda Equazione Cardinale}

La \textbf{seconda equazione cardinale} è l'equazione vettoriale che governa la dinamica rotazionale, e sta al momento angolare come il secondo principio di Newton sta alla quantità di moto. La ricaviamo in quattro casi distinti, combinando due alternative: il polo può essere \textit{fisso} o \textit{mobile}, e il sistema può essere un singolo \textit{punto materiale} oppure un \textit{sistema materiale}.

\subsection{Polo Fisso per un Punto Materiale}

Consideriamo un polo $O$ fisso, coincidente con il polo $O'$ rispetto a cui valutiamo il momento, per un singolo punto materiale. Deriviamo rispetto al tempo la definizione di momento angolare:
\begin{align}
    \left. \frac{d\vec{L}_O}{dt} \right|_{O} &= \left. \frac{d\vec{L}_O}{dt} \right|_{O'} = \frac{d}{dt}\left(\vec{r} \times m\vec{v}\right) \\
    &= m \, \frac{d\vec{r}}{dt} \times \vec{v} + m\,\vec{r} \times \frac{d\vec{v}}{dt} \\
    &= m\,(\vec{v} \times \vec{v}) + m\,(\vec{r} \times \vec{a})
\end{align}
Il primo termine è nullo, perché il prodotto vettoriale di un vettore con sé stesso si annulla. Nel secondo riconosciamo $m\vec{a} = \vec{F}$:
\begin{equation*}
    = \vec{r} \times m\vec{a} = \vec{r} \times \vec{F} = \vec{\tau}_O
\end{equation*}

\begin{teorema}{Seconda Equazione Cardinale --- Polo Fisso}
Per un punto materiale e un polo $O$ fisso vale, in \textbf{forma differenziale},
\begin{equation}
    \frac{d\vec{L}_O}{dt} = \vec{\tau}_O
\end{equation}
e, integrando fra due istanti $t_0$ e $t_1$, in \textbf{forma integrale},
\begin{equation}
    I_{\vec{\tau}_O} = \int_{t_0}^{t_1} \vec{\tau}_O(t') \, dt'
    = \int_{t_0}^{t_1} \frac{d\vec{L}_O}{dt'} \, dt'
    = \vec{L}_O(t_1) - \vec{L}_O(t_0) = \Delta\vec{L}_O
\end{equation}
L'impulso angolare del momento delle forze è dunque pari alla variazione del momento angolare.
\end{teorema}

\subsection{Polo Mobile per un Punto Materiale}

Supponiamo ora che $O$ resti fisso ma che il polo $O'$ rispetto a cui calcoliamo il momento sia \textbf{mobile}, con velocità $\vec{v}_{O'} = \frac{d\,\vec{OO'}}{dt} = -\frac{d\,\vec{O'O}}{dt}$. Partiamo dalla relazione fra i momenti rispetto ai due poli, $\vec{L}_{O'} = \vec{L}_O + \vec{O'O} \times m\vec{v}$, e deriviamo:
\begin{align}
    \left. \frac{d\vec{L}_{O'}}{dt} \right|_{O'}
    &= \frac{d}{dt}\left( \vec{L}_O + \vec{O'O} \times m\vec{v} \right) \\
    &= \frac{d\vec{L}_O}{dt} + \frac{d(\vec{O'O})}{dt} \times m\vec{v} + \vec{O'O} \times m\frac{d\vec{v}}{dt} \\
    &= \vec{\tau}_O + (-\vec{v}_{O'}) \times m\vec{v} + \vec{O'O} \times m\vec{a}
\end{align}
Riscriviamo ora $\vec{\tau}_O = \vec{OP} \times \vec{F}$ e raccogliamo i due termini che contengono la forza, ricomponendo il braccio rispetto al nuovo polo:
\begin{align}
    \left. \frac{d\vec{L}_{O'}}{dt} \right|_{O'}
    &= \big(\vec{OP} \times \vec{F} + \vec{O'O} \times \vec{F}\big) - \vec{v}_{O'} \times m\vec{v} \\
    &= (\vec{O'O} + \vec{OP}) \times \vec{F} - \vec{v}_{O'} \times m\vec{v} \\
    &= \vec{O'P} \times \vec{F} - \vec{v}_{O'} \times m\vec{v}
    = \vec{\tau}_{O'} - \vec{v}_{O'} \times m\vec{v}
\end{align}

\begin{teorema}{Seconda Equazione Cardinale --- Polo Mobile}
Per un punto materiale e un polo $O'$ in moto con velocità $\vec{v}_{O'}$ vale
\begin{equation}
    \frac{d\vec{L}_{O'}}{dt} + \vec{v}_{O'} \times m\vec{v} = \vec{\tau}_{O'}
\end{equation}
Rispetto al caso di polo fisso compare il termine aggiuntivo $\vec{v}_{O'} \times m\vec{v}$, che si annulla quando il polo è fermo ($\vec{v}_{O'} = \vec{0}$) oppure quando esso si muove parallelamente al punto ($\vec{v}_{O'} \parallel \vec{v}$).
\end{teorema}

\begin{osservazione}{Il Centro di Massa è un Polo Mobile Privilegiato}
Il secondo caso di annullamento è quello che rende il centro di massa un polo così comodo: scegliendo $O' \equiv CM$ si ha $\vec{v}_{O'} = \vec{V}_{CM}$ e, per un sistema, $\sum_i m_i\vec{v}_i = M\vec{V}_{CM}$, sicché il termine aggiuntivo diventa $\vec{V}_{CM} \times M\vec{V}_{CM} = \vec{0}$. L'equazione cardinale riprende allora la forma semplice valida per i poli fissi, anche se il centro di massa è in moto.
\end{osservazione}

\subsection{Polo Fisso per un Sistema Materiale}

Passiamo ora a un sistema di $N$ punti materiali, mantenendo il polo $O \equiv O'$ fisso. Il momento angolare totale è $\vec{L}_O^{\,\text{tot}} = \sum_i \vec{r}_i \times m_i \vec{v}_i$, e derivandolo termine a termine:
\begin{align}
    \left. \frac{d\vec{L}_O^{\,\text{tot}}}{dt} \right|_{O}
    &= \frac{d}{dt}\left( \sum_i \vec{r}_i \times m_i \vec{v}_i \right)
     = \sum_i \frac{d\vec{L}_{Oi}}{dt} \\
    &= \sum_i \vec{\tau}_{Oi}
     = \sum_i \vec{\tau}_{Oi}^{\,\text{INT}} + \sum_i \vec{\tau}_{Oi}^{\,\text{EXT}}
\end{align}
avendo applicato a ciascun punto il risultato del primo caso.

Il punto cruciale è che la sommatoria dei momenti delle forze \emph{interne} si annulla: per il terzo principio esse sono a due a due uguali e opposte, e agiscono lungo la congiungente dei due punti, sicché ogni coppia ha momento risultante nullo. Quindi $\sum_i \vec{\tau}_{Oi}^{\,\text{INT}} = \vec{0}$.

\begin{teorema}{Seconda Equazione Cardinale per i Sistemi}
Per un sistema materiale e un polo fisso $O$ vale, in \textbf{forma differenziale},
\begin{equation}
    \frac{d\vec{L}_O^{\,\text{tot}}}{dt} = \sum_i \vec{\tau}_{Oi}^{\,\text{EXT}}
\end{equation}
e, in \textbf{forma integrale},
\begin{equation}
    I^{\,\text{tot}}_{\vec{\tau}_O^{\,\text{EXT}}}
    = \int_{t_0}^{t_1} \sum_i \vec{\tau}_{Oi}^{\,\text{EXT}}(t') \, dt'
    = \Delta\vec{L}_O^{\,\text{tot}}
\end{equation}
Soltanto le forze \textbf{esterne} possono modificare il momento angolare complessivo di un sistema.
\end{teorema}

\subsection{Polo Mobile per un Sistema Materiale}

L'ultimo caso si ottiene sommando su tutte le particelle la relazione trovata per il punto materiale con polo mobile, $\vec{\tau}_{O'i} = \frac{d\vec{L}_{O'i}}{dt} + \vec{v}_{O'} \times m_i\vec{v}_i$:
\begin{equation}
    \sum_i \vec{\tau}_{O'i}^{\,\text{EXT}}
    = \frac{d\vec{L}_{O'}^{\,\text{tot}}}{dt} + \sum_i \vec{v}_{O'} \times m_i \vec{v}_i
\end{equation}
da cui l'espressione conclusiva
\begin{equation}
    \frac{d\vec{L}_{O'}^{\,\text{tot}}}{dt} + \vec{v}_{O'} \times M\vec{V}_{CM}
    = \sum_i \vec{\tau}_{O'i}^{\,\text{EXT}}
\end{equation}
avendo riconosciuto $\sum_i m_i\vec{v}_i = M\vec{V}_{CM}$, la quantità di moto totale del sistema.

\subsection{Conservazione del Momento Angolare}

Dalle equazioni cardinali discende immediatamente il teorema di conservazione, che è di gran lunga il loro impiego più frequente.

\begin{teorema}{Conservazione del Momento Angolare}
Se il sistema è \textbf{isolato}, o più in generale se la risultante dei momenti delle forze esterne è nulla o trascurabile,
\begin{equation*}
    \sum_i \vec{\tau}_{Oi}^{\,\text{EXT}} \approx \vec{0}
    \qquad \Longrightarrow \qquad
    \Delta\vec{L}_O^{\,\text{tot}} = \vec{0}
\end{equation*}
allora il momento angolare totale assume lo stesso valore in istanti diversi,
\begin{equation}
    \vec{L}_O^{\,\text{tot}}(t_1) = \vec{L}_O^{\,\text{tot}}(t_0)
\end{equation}
e costituisce una \textbf{costante del moto}.
\end{teorema}

\begin{attenzione}{Conservazione del Momento Angolare e della Quantità di Moto sono Indipendenti}
Le due leggi rispondono a condizioni diverse: $\vec{P}_{\text{tot}}$ si conserva se è nulla la \emph{risultante} delle forze esterne, $\vec{L}_O^{\,\text{tot}}$ se è nullo il loro \emph{momento risultante}. Una coppia di forze uguali e opposte applicate in punti diversi ha risultante nulla ma momento non nullo: conserva la quantità di moto e non il momento angolare. Viceversa, una forza centrale conserva il momento angolare rispetto al centro ma non la quantità di moto.
\end{attenzione}

\begin{osservazione}{Identità del Doppio Prodotto Vettoriale}
Nei calcoli che seguono comparirà ripetutamente il prodotto vettoriale triplo; conviene richiamarne l'identità fondamentale (regola \emph{BAC--CAB}):
\begin{equation*}
    \vec{A} \times (\vec{B} \times \vec{C})
    = \vec{B}\,(\vec{A} \cdot \vec{C}) - \vec{C}\,(\vec{A} \cdot \vec{B})
\end{equation*}
\end{osservazione}

\section{Teorema di Koenig per il Momento Angolare}

Come già per l'energia cinetica, anche il momento angolare di un sistema ammette una scomposizione che separa il moto d'insieme dal moto interno. Indichiamo con l'apice le grandezze riferite al sistema del centro di massa:
\begin{equation}
    \vec{r}_i = \vec{r}_i\,' + \vec{R}_{CM}
    \qquad \text{e} \qquad
    \vec{v}_i = \vec{v}_i\,' + \vec{V}_{CM}
\end{equation}

Calcoliamo il momento angolare totale rispetto al polo fisso $O$ e sviluppiamo il prodotto:
\begin{align}
    \vec{L}_O &= \sum_i \vec{r}_i \times m_i \vec{v}_i = \sum_i m_i (\vec{r}_i \times \vec{v}_i) \\
    &= \sum_i m_i \left[ (\vec{r}_i\,' + \vec{R}_{CM}) \times (\vec{v}_i\,' + \vec{V}_{CM}) \right] \\
    &= \sum_i m_i \left( \vec{r}_i\,' \times \vec{v}_i\,'
       + \vec{r}_i\,' \times \vec{V}_{CM}
       + \vec{R}_{CM} \times \vec{v}_i\,'
       + \vec{R}_{CM} \times \vec{V}_{CM} \right)
\end{align}

I due termini misti si annullano grazie alla definizione stessa di centro di massa, che impone $\sum_i m_i \vec{r}_i\,' = \vec{0}$ e, derivando, $\sum_i m_i \vec{v}_i\,' = \vec{0}$:
\begin{align}
    \sum_i m_i (\vec{r}_i\,' \times \vec{V}_{CM})
    &= \Big(\underbrace{\textstyle\sum_i m_i \vec{r}_i\,'}_{=\,\vec{0}}\Big) \times \vec{V}_{CM} = \vec{0} \\
    \sum_i m_i (\vec{R}_{CM} \times \vec{v}_i\,')
    &= \vec{R}_{CM} \times \Big(\underbrace{\textstyle\sum_i m_i \vec{v}_i\,'}_{=\,\vec{0}}\Big) = \vec{0}
\end{align}

Restano soltanto il primo e il quarto termine, e quest'ultimo si compatta osservando che $\sum_i m_i = M$:
\begin{equation}
    \vec{L}_O = \underbrace{\sum_i \vec{r}_i\,' \times m_i \vec{v}_i\,'}_{\vec{L}_{O'}}
    + \underbrace{M \,(\vec{R}_{CM} \times \vec{V}_{CM})}_{\vec{L}_{CM}}
\end{equation}

\begin{teorema}{Teorema di Koenig per il Momento Angolare}
Il momento angolare totale di un sistema rispetto a un polo fisso $O$ si decompone nella somma di due contributi indipendenti:
\begin{equation}
    \vec{L}_O = \vec{L}_{O'} + \vec{L}_{CM}
\end{equation}
dove
\begin{itemize}[leftmargin=1.4em]
    \item $\vec{L}_O$ è il momento angolare totale del sistema rispetto al polo fisso $O$;
    \item $\vec{L}_{O'}$ è il momento angolare \textbf{intrinseco} (o di spin), dovuto al moto delle particelle rispetto al centro di massa;
    \item $\vec{L}_{CM} = M(\vec{R}_{CM} \times \vec{V}_{CM})$ è il momento angolare \textbf{orbitale}, quello che avrebbe un unico punto materiale di massa $M$ posto nel centro di massa.
\end{itemize}
\end{teorema}

\begin{osservazione}{Orbitale e Intrinseco}
La distinzione ricalca quella fra il moto di rivoluzione della Terra attorno al Sole --- che contribuisce a $\vec{L}_{CM}$ --- e il suo moto di rotazione attorno al proprio asse, che contribuisce a $\vec{L}_{O'}$. I due termini sono indipendenti: si possono modificare l'uno senza toccare l'altro.
\end{osservazione}

\subsection{Caso del Sistema a Due Punti Materiali}

Applichiamo il teorema al caso più importante, quello di un sistema isolato di due soli punti materiali. Accanto alle coordinate del centro di massa introduciamo le coordinate \textbf{relative}, che descrivono la posizione di un punto rispetto all'altro:
\begin{equation}
    \vec{R}_{CM} = \frac{m_1 \vec{r}_1 + m_2 \vec{r}_2}{m_1 + m_2}
    \, ; \qquad
    \begin{cases}
        \vec{r}_r = \vec{r}_2 - \vec{r}_1 \\
        \vec{v}_r = \vec{v}_2 - \vec{v}_1
    \end{cases}
\end{equation}

Invertendo queste relazioni si esprimono le posizioni assolute in funzione delle due nuove variabili:
\begin{equation}
    \begin{cases}
        \vec{r}_1 = \vec{R}_{CM} - \dfrac{m_2}{m_1 + m_2}\,\vec{r}_r \\[8pt]
        \vec{r}_2 = \vec{R}_{CM} + \dfrac{m_1}{m_1 + m_2}\,\vec{r}_r
    \end{cases}
\end{equation}
e, sottraendo $\vec{R}_{CM}$, si ottengono le posizioni e le velocità nel riferimento del centro di massa:
\begin{equation}
    \begin{cases}
        \vec{r}_1\,' = - \dfrac{m_2}{m_1 + m_2}\,\vec{r}_r \\[8pt]
        \vec{r}_2\,' = + \dfrac{m_1}{m_1 + m_2}\,\vec{r}_r
    \end{cases}
    \qquad \Longrightarrow \qquad
    \begin{cases}
        \vec{v}_1\,' = - \dfrac{m_2}{m_1 + m_2}\,\vec{v}_r \\[8pt]
        \vec{v}_2\,' = + \dfrac{m_1}{m_1 + m_2}\,\vec{v}_r
    \end{cases}
\end{equation}

Sostituiamo nell'espressione del momento angolare intrinseco:
\begin{align}
    \vec{L}_{O'} &= m_1 (\vec{r}_1\,' \times \vec{v}_1\,') + m_2 (\vec{r}_2\,' \times \vec{v}_2\,') \\
    &= m_1 \left( \frac{m_2}{m_1+m_2} \right)^{\!2} (\vec{r}_r \times \vec{v}_r)
     + m_2 \left( \frac{m_1}{m_1+m_2} \right)^{\!2} (\vec{r}_r \times \vec{v}_r) \\
    &= \frac{m_1 m_2^2 + m_2 m_1^2}{(m_1+m_2)^2} \, (\vec{r}_r \times \vec{v}_r) \\
    &= \frac{m_1 m_2 \,\cancel{(m_1 + m_2)}}{(m_1+m_2)^{\cancel{2}}} \, (\vec{r}_r \times \vec{v}_r)
     = \frac{m_1 m_2}{m_1+m_2} \, (\vec{r}_r \times \vec{v}_r)
\end{align}
(entrambi i coefficienti compaiono elevati al quadrato, dunque il segno meno di $\vec{r}_1\,'$ e $\vec{v}_1\,'$ scompare e i due contributi si sommano).

\begin{definizione}{Massa Ridotta}
Il fattore scalare che compare a moltiplicare il prodotto vettoriale prende il nome di \textbf{massa ridotta} del sistema:
\begin{equation}
    \mu = \frac{m_1 m_2}{m_1+m_2}
\end{equation}
Essa è sempre minore della più piccola delle due masse, e tende a quest'ultima quando l'altra diventa molto grande.
\end{definizione}

Il momento angolare intrinseco assume dunque la forma compatta
\begin{equation}
    \vec{L}_{O'} = \mu \,(\vec{r}_r \times \vec{v}_r)
\end{equation}
che è formalmente identica al momento angolare di \emph{un solo} punto materiale, di massa $\mu$, posizione $\vec{r}_r$ e velocità $\vec{v}_r$. Il problema a due corpi si è ridotto a un problema a un corpo: chiameremo \textbf{particella equivalente} questo punto fittizio di parametri $(\mu, \vec{r}_r, \vec{v}_r)$.

\subsection{Equivalenza del Sistema a Due Corpi}

La riduzione appena ottenuta è un'identità algebrica; resta da verificare che essa si estenda anche alla \emph{dinamica}, cioè che la particella equivalente obbedisca effettivamente a un'equazione del moto del tipo previsto. Partiamo dalla seconda equazione cardinale per il sistema, scomponendo le posizioni rispetto al centro di massa:
\begin{align}
    \frac{d\vec{L}_O^{\,\text{tot}}}{dt}
    &= \sum_i \vec{\tau}_{Oi}^{\,\text{EXT}} = \sum_i (\vec{r}_i \times \vec{F}_i^{\,\text{EXT}})
     = \sum_i (\vec{r}_i\,' + \vec{R}_{CM}) \times \vec{F}_i^{\,\text{EXT}} \\
    &= \sum_i (\vec{r}_i\,' \times \vec{F}_i^{\,\text{EXT}}) + \sum_i (\vec{R}_{CM} \times \vec{F}_i^{\,\text{EXT}}) \\
    &= \sum_i \vec{\tau}_{O'i}^{\,\text{EXT}} + \vec{R}_{CM} \times \vec{F}_{\text{TOT}}^{\,\text{EXT}}
\end{align}
che riproduce la scomposizione di Koenig anche a livello di derivate:
\begin{equation*}
    \frac{d\vec{L}_O^{\,\text{tot}}}{dt}
    = \frac{d\vec{L}_{O'}^{\,\text{tot}}}{dt} + \vec{R}_{CM} \times \frac{d\vec{P}^{\,\text{TOT}}}{dt}
\end{equation*}

Calcoliamo ora esplicitamente il termine intrinseco per il sistema a due corpi:
\begin{align}
    \frac{d\vec{L}_{O'}}{dt}
    &= \frac{d}{dt}\left( m_1 (\vec{r}_1\,' \times \vec{v}_1\,') + m_2 (\vec{r}_2\,' \times \vec{v}_2\,') \right)
\end{align}
I termini $\vec{v}_i\,' \times \vec{v}_i\,'$ si annullano, e ricordando che $\vec{v}_i\,' = \vec{v}_i - \vec{V}_{CM}$ resta
\begin{align}
    &= m_1 \left( \vec{r}_1\,' \times \frac{d\vec{v}_1}{dt} \right)
     + m_2 \left( \vec{r}_2\,' \times \frac{d\vec{v}_2}{dt} \right)
     - \big(\underbrace{m_1 \vec{r}_1\,' + m_2 \vec{r}_2\,'}_{=\,\vec{0}}\big) \times \frac{d\vec{V}_{CM}}{dt}
\end{align}
dove l'ultimo termine si annulla per la definizione di centro di massa. Traducendo le accelerazioni in forze:
\begin{align}
    \frac{d\vec{L}_{O'}}{dt}
    &= \vec{r}_1\,' \times \vec{F}_1 + \vec{r}_2\,' \times \vec{F}_2 \\
    &= \vec{r}_1\,' \times (\vec{F}_1^{\,\text{EXT}} + \vec{F}_{21})
     + \vec{r}_2\,' \times (\vec{F}_2^{\,\text{EXT}} + \vec{F}_{12})
\end{align}
Il terzo principio impone $\vec{F}_{12} = -\vec{F}_{21}$, e raccogliendo:
\begin{align}
    &= \big(\vec{r}_1\,' \times \vec{F}_1^{\,\text{EXT}} + \vec{r}_2\,' \times \vec{F}_2^{\,\text{EXT}}\big)
       + (\vec{r}_1\,' - \vec{r}_2\,') \times \vec{F}_{21} \\
    &= \vec{\tau}_{O'}^{\,\text{EXT}} + (\vec{r}_1\,' - \vec{r}_2\,') \times \vec{F}_{21}
\end{align}

Resta da valutare il termine dovuto alle forze interne. Esprimiamolo nella coordinata relativa:
\begin{equation*}
    \vec{r}_1\,' - \vec{r}_2\,'
    = (\vec{r}_1 - \vec{R}_{CM}) - (\vec{r}_2 - \vec{R}_{CM})
    = \vec{r}_1 - \vec{r}_2 = -\,\vec{r}_r
\end{equation*}
Il contributo interno vale dunque $-\,\vec{r}_r \times \vec{F}_{21}$. Ma la forza che i due corpi si scambiano è diretta lungo la loro congiungente --- è cioè una forza centrale rispetto all'altro corpo --- e quindi $\vec{F}_{21} \parallel \vec{r}_r$. Il prodotto vettoriale fra vettori paralleli è nullo:
\begin{equation}
    -\,\vec{r}_r \times \vec{F}_{21} = \vec{0}
\end{equation}

\begin{teorema}{Riduzione del Problema a Due Corpi}
Un sistema isolato di due corpi interagenti equivale in tutto e per tutto a una singola \textbf{particella equivalente} di parametri $(\mu, \vec{r}_r, \vec{v}_r)$, soggetta alla sola forza centrale diretta lungo la congiungente $\vec{r}_r$.
\end{teorema}

\begin{osservazione}{Perché la Riduzione è Così Utile}
Il teorema trasforma un problema a sei gradi di libertà (due punti nello spazio) in due problemi indipendenti e molto più semplici: il moto rettilineo uniforme del centro di massa, che è banale, e il moto di un singolo punto in un campo centrale, per il quale disponiamo già delle due costanti $E$ e $\vec{L}$. È su quest'ultimo che si concentra tutto il resto del capitolo.
\end{osservazione}

\section{Energia e Potenziale Efficace}

Studiamo ora l'energia del sistema dal punto di vista della particella equivalente. Il suo momento angolare, che sappiamo costante, si scrive in coordinate polari come
\begin{equation}
    \vec{L}_{O'} = \mu \,(\vec{r}_r \times \vec{v}_r)
    = \mu \, \big(r_r \hat{r} \times (\dot{r}_r \hat{r} + r_r \dot{\phi} \hat{\phi})\big)
    = \mu \, r_r^2 \, \omega \, \hat{z}
\end{equation}

L'energia totale è la somma di energia cinetica ed energia potenziale; esplicitando la velocità in coordinate polari:
\begin{align}
    E_{\text{tot}} &= K + U = \frac{1}{2} \mu \, \vec{v}_r^{\;2} + U(r) \\
    &= \frac{1}{2} \mu \left( \dot{r}_r^{\,2} + r_r^2 \dot{\phi}^2 \right) + U(r)
\end{align}

Il passaggio decisivo consiste nell'eliminare $\dot{\phi}$ servendosi della conservazione del momento angolare. Da $L = \mu r_r^2 \dot{\phi}$ segue $\dot{\phi} = \dfrac{L}{\mu r_r^2}$, e sostituendo:
\begin{equation}
    E_{\text{tot}} = \frac{1}{2} \mu \dot{r}_r^{\,2} + \frac{1}{2}\,\frac{L^2}{\mu r_r^2} + U(r)
\end{equation}
L'energia dipende ora dalla sola coordinata radiale: il problema bidimensionale è diventato \textbf{unidimensionale}.

\subsection{Definizione di Potenziale Efficace}

Raggruppiamo tutti i termini che dipendono esclusivamente da $r$:
\begin{equation}
    E_{\text{tot}} = \underbrace{\frac{1}{2} \mu \dot{r}_r^{\,2}}_{K_{\text{rad}}}
    + \underbrace{\frac{1}{2}\,\frac{L^2}{\mu r_r^2} + U(r)}_{U_{\text{eff}}(r)}
\end{equation}
dove si riconoscono:
\begin{itemize}
    \item $\frac{1}{2}\mu \dot{r}_r^{\,2}$: l'\textbf{energia cinetica radiale};
    \item $\frac{L^2}{2\mu r^2}$: il \textbf{potenziale centrifugo}, che non è una vera energia potenziale ma il residuo dell'energia cinetica di rotazione, riscritto come funzione di $r$;
    \item $U(r)$: il potenziale dell'interazione centrale vera e propria.
\end{itemize}

\begin{definizione}{Potenziale Efficace}
Si definisce \textbf{potenziale efficace} la funzione
\begin{equation}
    U_{\text{eff}}(r) = U(r) + \frac{L^2}{2\mu r^2}
\end{equation}
Grazie a essa il moto radiale della particella equivalente si tratta esattamente come un moto unidimensionale in un campo di energia potenziale $U_{\text{eff}}$.
\end{definizione}

\begin{attenzione}{Il Potenziale Centrifugo Dipende dalle Condizioni Iniziali}
A differenza di $U(r)$, che è fissato una volta per tutte dalla natura dell'interazione, il termine centrifugo contiene $L$, che dipende da \emph{come} il moto è stato avviato. Cambiando il momento angolare iniziale cambia la forma di $U_{\text{eff}}$, e con essa il tipo di orbita: lo stesso campo di forze produce traiettorie qualitativamente diverse a seconda di $L$.
\end{attenzione}

\subsection{Condizione di Esistenza del Moto}

Perché il moto sia fisicamente possibile l'energia cinetica radiale deve essere non negativa. Questa richiesta si traduce in un vincolo fra l'energia totale e il potenziale efficace:
\begin{equation}
    \frac{1}{2} \mu \dot{r}^2 \geq 0
    \quad \Longrightarrow \quad E - U_{\text{eff}}(r) \geq 0
    \quad \Longrightarrow \quad E \geq U_{\text{eff}}(r)
\end{equation}
La condizione delimita i valori di $r$ che la particella può effettivamente assumere, per una data energia $E$.

\subsection{Regioni Accessibili e Punti di Inversione}

Il moto è quindi consentito solo dove la curva $U_{\text{eff}}(r)$ sta al di sotto del livello $E$; le regioni in cui ciò accade si dicono \textbf{accessibili}, le altre proibite. Il segno della velocità radiale distingue le due fasi del moto:
\begin{itemize}
    \item $\dot{r} > 0$: la particella si \textbf{allontana} dal centro di forza;
    \item $\dot{r} < 0$: la particella si \textbf{avvicina} al centro di forza;
    \item $\dot{r} = 0$: la particella ha esaurito l'energia disponibile per il moto radiale ed è costretta a invertire il verso. Il punto corrispondente si dice \textbf{punto di inversione} (o punto nodale).
\end{itemize}

Nei punti di inversione la velocità è puramente tangenziale, poiché la componente radiale si annulla:
\begin{equation}
    \vec{v} = \cancelto{0}{\dot{r}\hat{r}} + r\dot{\phi}\hat{\phi} = r\omega\hat{\phi} = \vec{v}_{\text{tang}}
\end{equation}

\begin{figure}[ht]
    \centering
    \begin{tikzpicture}[x=1.55cm, y=1.15cm, >=stealth, font=\small]
        % Assi
        \draw[->, thick, mypen] (0,-2.6) -- (0,3.6) node[above] {$U_{\text{eff}}(r)$};
        \draw[->, thick, mypen] (-0.2,0) -- (6.3,0) node[right] {$r$};

        % Livelli energetici
        \draw[dashed, boxoss]  (0,2.5)  node[left, boxoss]  {$E_1$} -- (5.8,2.5);
        \draw[dashed, boxatt]  (0,0.8)  node[left, boxatt]  {$E_2$} -- (5.8,0.8);
        \draw[dashed, boxese]  (0,-0.5) node[left, boxese]  {$E_3$} -- (5.8,-0.5);
        \draw[dashed, boxnota] (0,-1.2) node[left, boxnota] {$E_4$} -- (3.5,-1.2);
        \draw[dashed, notered] (0,-1.8) node[left, notered] {$E_5$} -- (2.6,-1.8);

        % Curva del potenziale efficace
        \draw[ultra thick, mypen, smooth, tension=0.7] plot coordinates {
            (0.35,3.4) (0.6,1.0) (1.5,-1.8) (2.8,0.8) (4.5,0.2) (5.8,0.05)
        };
        \node[mypen, right] at (5.0,0.62) {$U_{\text{eff}}(r)$};

        % Punti di inversione
        \fill[boxoss]  (0.42,2.5)  circle (2pt) node[above right, boxoss, xshift=-2pt]  {$r_0$};
        \fill[boxatt]  (0.65,0.8)  circle (2pt) node[below left, boxatt, xshift=4pt]    {$r_1$};
        \fill[boxatt]  (2.8,0.8)   circle (2pt) node[above, boxatt]                     {$r_7$};
        \fill[boxese]  (0.90,-0.5) circle (2pt) node[below left, boxese, xshift=4pt]    {$r_2$};
        \fill[boxese]  (2.13,-0.5) circle (2pt) node[below right, boxese, xshift=-3pt]  {$r_6$};
        \fill[boxese]  (3.80,-0.5) circle (2pt) node[above right, boxese, xshift=-3pt]  {$r_8$};
        \fill[boxnota] (1.10,-1.2) circle (2pt) node[below left, boxnota, xshift=5pt]   {$r_3$};
        \fill[boxnota] (1.85,-1.2) circle (2pt) node[below right, boxnota, xshift=-4pt] {$r_5$};
        \fill[notered] (1.50,-1.8) circle (2pt) node[below, notered]                    {$r_4$};
    \end{tikzpicture}
    \caption{Andamento qualitativo del potenziale efficace e classificazione delle orbite. Ogni livello orizzontale rappresenta un valore dell'energia totale; le sue intersezioni con la curva sono i punti di inversione $r_i$, che delimitano le regioni accessibili al moto.}
    \label{fig:potenziale_efficace}
\end{figure}

Il tipo di orbita dipende dunque interamente dalle \textbf{condizioni iniziali}, riassunte nella coppia $(L_0, E_0)$. Scorrendo la figura dall'alto verso il basso:

\begin{itemize}
    \item $\mathbf{E_0 = E_1}$: il moto è accessibile per $r \geq r_0$. La particella giunge da grande distanza, raggiunge la minima distanza $r_0$ e torna verso l'infinito: si tratta di un \textbf{moto aperto} (illimitato).
    \item $\mathbf{E_0 = E_2}$: il moto è accessibile per $r \geq r_1$. Oltre alla traiettoria aperta di avvicinamento e allontanamento, in $r = r_7$ --- dove il livello è tangente al \emph{massimo} della curva --- esiste un'orbita circolare, ma \textbf{instabile}: una perturbazione infinitesima allontana definitivamente la particella da essa.
    \item $\mathbf{E_0 = E_3}$: la curva presenta ora due regioni accessibili separate. La particella può restare intrappolata nella buca, $r_2 \leq r \leq r_6$, oppure muoversi liberamente all'esterno per $r \geq r_8$. Quale dei due destini si realizzi dipende da dove essa si trovi inizialmente: le due regioni non comunicano.
    \item $\mathbf{E_0 = E_4}$: il moto è accessibile solo per $r_3 \leq r \leq r_5$. La particella oscilla confinata fra due \textbf{absidi}, alternando avvicinamenti e allontanamenti: il \textbf{moto è limitato}.
    \item $\mathbf{E_0 = E_5}$: il livello è tangente al \emph{minimo} della curva, dunque $K_{\text{rad}} = 0$ in permanenza e l'unico valore possibile è $r = r_4$. Il raggio non varia mai: l'orbita è \textbf{circolare e stabile}, e corrisponde alla minima energia compatibile con il momento angolare assegnato.
    \item $\mathbf{E_0 < E_5}$: si avrebbe $K_{\text{rad}} < 0$, che è impossibile. Nessun moto è consentito.
\end{itemize}

\begin{osservazione}{Stabilità e Concavità}
La differenza fra l'orbita circolare in $r_7$ e quella in $r_4$ non sta nell'energia ma nella \emph{concavità} della curva: entrambe annullano $dU_{\text{eff}}/dr$, ma solo il minimo ($d^2U_{\text{eff}}/dr^2 > 0$) richiama la particella verso l'orbita dopo una perturbazione. Il massimo la respinge, ed è per questo instabile.
\end{osservazione}

\subsection{Orbite Circolari nei Sistemi a Due Corpi}

Il caso del minimo di $U_{\text{eff}}$ merita un esame a parte, perché descrive le orbite circolari --- la configurazione di gran lunga più comune nei sistemi astronomici. Per una singola particella la condizione è semplicemente $r = \text{costante}$; vediamo che cosa significhi per un sistema a due corpi.

Nel sistema isolato la particella equivalente ha raggio costante, e dunque resta costante la \textbf{distanza relativa} fra i due corpi:
\begin{equation}
    r_r = |\vec{r}_2 - \vec{r}_1| = |\vec{r}_2\,' - \vec{r}_1\,'| = \text{costante}
\end{equation}

Richiamando le posizioni rispetto al centro di massa già ricavate,
\begin{equation}
    \begin{cases}
        \vec{r}_1\,' = - \dfrac{m_2}{m_1+m_2}\,\vec{r}_r \\[8pt]
        \vec{r}_2\,' = + \dfrac{m_1}{m_1+m_2}\,\vec{r}_r
    \end{cases}
\end{equation}
si vede subito che, essendo $|\vec{r}_r|$ costante, sono costanti anche $|\vec{r}_1\,'|$ e $|\vec{r}_2\,'|$. Da questa sola osservazione discendono tre conclusioni.

\begin{enumerate}
    \item \textbf{Entrambi i corpi percorrono un'orbita circolare} attorno al centro di massa, con raggi $|\vec{r}_1\,'|$ e $|\vec{r}_2\,'|$ inversamente proporzionali alle rispettive masse: il corpo più pesante descrive il cerchio più piccolo.
    \item \textbf{La velocità angolare è la stessa per entrambi} ed è costante. I due corpi restano infatti sempre diametralmente opposti rispetto al centro di massa --- i vettori $\vec{r}_1\,'$ e $\vec{r}_2\,'$ sono antiparalleli per costruzione --- e quindi spazzano lo stesso angolo nello stesso tempo.
    \item \textbf{Il centro di forza di ciascun corpo è l'altro corpo}, non il centro di massa. Quest'ultimo è solo il punto attorno al quale le due orbite appaiono descritte: nessuna massa vi è collocata e nessuna forza vi è applicata.
\end{enumerate}

\begin{figure}[ht]
    \centering
    \begin{tikzpicture}[>=stealth, scale=1.5, font=\small]
        \def\rA{1.15}   % raggio del corpo pesante (orbita interna)
        \def\rB{2.50}   % raggio del corpo leggero (orbita esterna)
        \def\angT{250}  % posizione di m_1 all'istante t
        \def\angP{210}  % posizione di m_1 all'istante t'

        % Orbite
        \draw[dashed, thick, mypen!40] (0,0) circle (\rA);
        \draw[dashed, thick, mypen!40] (0,0) circle (\rB);

        % Verso di rotazione
        \draw[->, thick, mypen!65] (100:\rB+0.30) arc (100:140:\rB+0.30);
        \node[mypen!80] at (120:\rB+0.58) {$\vec{\omega}$};

        % ---- istante t' (disegnato prima, in secondo piano) ----
        \coordinate (Ap) at (\angP:\rA);
        \coordinate (Bp) at ({\angP+180}:\rB);
        \draw[thick, mypen!40, dashed] (Ap) -- (Bp);
        \filldraw[notered!18, draw=notered!55, thick] (Ap) circle (0.16);
        \filldraw[boxoss!18, draw=boxoss!55, thick] (Bp) circle (0.12);
        \node[notered!75] at (\angP:\rA+0.48) {$m_1(t')$};
        \node[boxoss!75] at ({\angP+180}:\rB+0.42) {$m_2(t')$};

        % ---- istante t ----
        \coordinate (A) at (\angT:\rA);
        \coordinate (B) at ({\angT+180}:\rB);
        \draw[very thick, mypen] (A) -- (B);
        \filldraw[notered!45, draw=notered, very thick] (A) circle (0.16);
        \filldraw[boxoss!45, draw=boxoss, very thick] (B) circle (0.12);
        \node[notered] at (\angT:\rA+0.48) {$m_1(t)$};
        \node[boxoss] at ({\angT+180}:\rB+0.40) {$m_2(t)$};

        % Raggi rispetto al centro di massa
        \draw[->, very thick, notered] (0,0) -- (A)
            node[pos=0.62, anchor=west, notered, xshift=2pt] {$\vec{r}_1\,'$};
        \draw[->, very thick, boxoss] (0,0) -- (B)
            node[pos=0.55, anchor=east, boxoss, xshift=-2pt] {$\vec{r}_2\,'$};

        % Angoli spazzati, uguali sui due lati
        \draw[thick, boxese] (\angT:0.58) arc (\angT:\angP:0.58);
        \node[boxese] at ({(\angT+\angP)/2}:0.82) {$\theta$};
        \draw[thick, boxese] ({\angT+180}:1.62) arc ({\angT+180}:{\angP+180}:1.62);
        \node[boxese] at ({(\angT+\angP)/2+180}:1.86) {$\theta$};

        % Centro di massa
        \filldraw[boxese] (0,0) circle (2.0pt);
        \node[boxese, anchor=south west] at (0.10,0.10) {$CM$};
    \end{tikzpicture}
    \caption{Orbite circolari in un sistema isolato a due corpi, ritratte in due istanti successivi $t$ e $t'$. I due corpi restano costantemente allineati con il centro di massa e su lati opposti rispetto a esso, e spazzano nello stesso tempo il medesimo angolo $\theta$: la velocità angolare è comune. Il corpo più pesante ($m_1$) percorre l'orbita di raggio minore.}
    \label{fig:orbite_circolari_dettagliate}
\end{figure}

\subsection{Orbite Circolari e Punti Stazionari del Potenziale Efficace}

Verifichiamo ora per via analitica quanto la figura mostra geometricamente, cercando i punti stazionari del potenziale efficace:
\begin{equation}
    U_{\text{eff}}(r) = \frac{L_0^2}{2\mu r^2} + U(r)
\end{equation}
Deriviamo rispetto al raggio:
\begin{align}
    \frac{dU_{\text{eff}}}{dr}
    &= \frac{dU}{dr} + \frac{L_0^2}{2\mu} \left( -\frac{2}{r^3} \right)
     = \frac{dU}{dr} - \frac{L_0^2}{\mu r^3} \\
    &= \frac{dU}{dr} - \frac{(\mu r^2 \omega)^2}{\mu r^3}
     = \frac{dU}{dr} - \frac{\mu^{\cancel{2}} r^{\cancel{4}} \omega^2}{\cancel{\mu} \, r^{\cancel{3}}}
     = \frac{dU}{dr} - \mu r \omega^2
\end{align}
avendo sostituito $L_0 = \mu r^2 \omega$. Nel caso del moto circolare il raggio si colloca in un punto stazionario, dove la derivata si annulla:
\begin{equation}
    \frac{dU_{\text{eff}}}{dr} = 0
    \quad \Longrightarrow \quad
    \frac{dU}{dr} - \mu r \omega^2 = 0
\end{equation}

Ricordiamo infine che una forza conservativa è l'opposto del gradiente del potenziale; proiettando lungo $\hat{r}$:
\begin{equation}
    -\frac{dU}{dr}\,\hat{r} = -\nabla U = \vec{F}
    \quad \Longrightarrow \quad
    \left(-\frac{dU}{dr}\hat{r}\right) \cdot \hat{r} = \vec{F} \cdot \hat{r} = F(r)
\end{equation}
cioè $-\dfrac{dU}{dr} = F(r)$. Sostituendo nella condizione di stazionarietà si ottiene
\begin{equation}
    -F(r) - \mu r \omega^2 = 0
    \quad \Longrightarrow \quad
    F(r) = - \mu r \omega^2
\end{equation}

\begin{osservazione}{Una Verifica di Coerenza}
Il risultato non è una nuova legge, ma la conferma che tutto il formalismo è consistente: $-\mu r\omega^2$ è esattamente la forza centripeta necessaria a mantenere la particella equivalente su una circonferenza di raggio $r$ alla velocità angolare $\omega$. Il segno negativo indica che la forza è diretta verso il centro, come dev'essere. Ritrovare per questa via un risultato elementare della cinematica circolare è la miglior verifica formale dei passaggi svolti.
\end{osservazione}

\subsection{Soluzione Formale per il Moto Radiale}

Concludiamo ricavando la legge oraria del moto radiale. Dalla conservazione dell'energia si esplicita la velocità radiale:
\begin{equation}
    \frac{1}{2}\mu \dot{r}^2 = E - U_{\text{eff}}(r)
    \quad \Longrightarrow \quad
    \dot{r} = \pm \sqrt{\frac{2}{\mu}\big(E - U_{\text{eff}}(r)\big)}
\end{equation}
Il doppio segno corrisponde alle due fasi del moto:
\begin{itemize}
    \item segno $(+)$, con $\dot{r} > 0$: fase di \textbf{allontanamento} dal centro di forza;
    \item segno $(-)$, con $\dot{r} < 0$: fase di \textbf{avvicinamento} al centro di forza.
\end{itemize}

Scrivendo $\dot{r} = dr/dt$ e separando le variabili:
\begin{equation}
    \frac{dr}{\sqrt{\dfrac{2}{\mu}\big(E - U_{\text{eff}}\big)}} = \pm \, dt
\end{equation}
Integrando fra l'istante iniziale $t_0$, in cui il raggio vale $r_0$, e l'istante generico $t$:
\begin{equation}
    \int_{r_0}^{r(t)} \frac{dr'}{\sqrt{\dfrac{2}{\mu}\big(E - U_{\text{eff}}(r')\big)}} = \pm \,(t - t_0)
\end{equation}

\begin{formulario}{Soluzione Formale del Moto Radiale}
\begin{equation}
    t = t_0 \pm \int_{r_0}^{r(t)} \frac{dr'}{\sqrt{\dfrac{2}{\mu}\big(E - U_{\text{eff}}(r')\big)}}
\end{equation}
\end{formulario}

\begin{osservazione}{Che Cosa Significa ``Formale''}
La soluzione è detta formale perché fornisce $t$ in funzione di $r$, cioè la funzione \emph{inversa} di quella cercata: per ottenere $r(t)$ occorre risolvere l'integrale --- il che è possibile in forma chiusa solo per pochi potenziali notevoli, fra cui quello kepleriano --- e poi invertire il risultato.

Ciò nonostante l'espressione è tutt'altro che inutile. Il doppio segno permette di ricostruire il moto a tratti, raccordando gli intervalli in cui $r(t)$ cresce ($\dot{r} > 0$) con quelli in cui decresce ($\dot{r} < 0$), separati dai punti di inversione. E anche senza risolverla si può usarla per calcolare quantità globali: valutata fra due punti di inversione consecutivi, essa fornisce direttamente il \textbf{periodo} dell'oscillazione radiale.
\end{osservazione}
